\chapter{A Full Practice Assessment}
\label{ch:exam}

This chapter is a practice assessment modeled on what candidates report: hardware multiple-choice and true/false questions, short C output questions, a few aptitude questions, and two coding problems. Take it under real conditions: no notes, a timer, and no checking answers until the end.

\begin{remember}
Suggested timing: 60 minutes for Parts A to C, then 30 minutes for Part D. If your assessment uses negative marking, skip questions you cannot narrow down to two options.
\end{remember}

\newcommand{\opts}[4]{\par\hspace*{1em}(a) #1\quad (b) #2\quad (c) #3\quad (d) #4}

\section{Part A: Hardware (multiple choice and true/false)}

\begin{enumerate}
  \item An NMOS pass transistor with its gate at $\VDD = 1.0$\,V and $\Vtn = 0.3$\,V passes a logic 1 to a capacitive load. The load settles at:
  \opts{1.0\,V}{0.3\,V}{0.7\,V}{0\,V}

  \item In a static CMOS gate, the network that connects the output to $\VDD$ is made of:
  \opts{NMOS in series}{PMOS transistors}{NMOS in parallel}{resistors}

  \item An inverter is built with its NMOS connected to $\VDD$ and its PMOS to ground (gates tied together as the input, sources tied as the output). For input 1, the output is approximately:
  \opts{0\,V}{$|\Vtp|$}{$\VDD$}{$\VDD - \Vtn$}

  \item The circuit in question 3 behaves as:
  \opts{an inverter with full swing}{an oscillator}{a non-inverting buffer with degraded levels}{an output stuck at 0}

  \item One advantage of a flip-flop over a latch is that a flip-flop:
  \opts{uses fewer transistors}{allows time borrowing}{has a lower clock load}{captures its input only at the clock edge, so data cannot race through}

  \item A 2:1 mux has $S = B$, $I_0 = A$, $I_1 = 1$. Its output is:
  \opts{$A + B$}{$AB$}{$A \oplus B$}{$\nt{A}$}

  \item The minimum number of 2:1 muxes needed to build XOR, if only $A$, $B$, 0 and 1 are available, is:
  \opts{1}{2}{3}{4}

  \item The clock frequency and supply voltage of a chip both double. Dynamic power becomes:
  \opts{$2\times$}{$4\times$}{$16\times$}{$8\times$}

  \item Supply voltage rises 10\% and frequency rises 20\%. Dynamic power rises by about:
  \opts{30\%}{45\%}{32\%}{21\%}

  \item As temperature increases, carrier mobility \underline{\hspace{1.5cm}} and metal wire resistance \underline{\hspace{1.5cm}}.
  \opts{increases; increases}{increases; decreases}{decreases; increases}{decreases; decreases}

  \item True or false: for the antenna rule, connecting a long wire to a \emph{smaller} gate reduces the risk of oxide damage.

  \item Which of these fixes an antenna violation?
  \opts{adding a reverse-biased diode near the gate}{making the gate smaller}{making the lower-layer wire longer}{raising $\VDD$}

  \item A path has a hold violation. Which fixes it?
  \opts{lowering the clock frequency}{raising the clock frequency}{removing buffers from the short path}{adding delay to the short path}

  \item $t_{cq} = 50$\,ps, $t_{pd} = 350$\,ps, $t_{su} = 100$\,ps, no skew. The maximum clock frequency is:
  \opts{1\,GHz}{2\,GHz}{2.5\,GHz}{2.86\,GHz}

  \item The capture flip-flop's clock arrives later than the launch flip-flop's. This:
  \opts{helps setup and hurts hold}{hurts setup and helps hold}{helps both}{affects neither}

  \item A T flip-flop with $T = 1$ is clocked at 25\,kHz. Its output frequency is:
  \opts{50\,kHz}{25\,kHz}{12.5\,kHz}{6.25\,kHz}

  \item A binary counter's clock period is 2\,ns and its most significant bit has a period of 128\,ns. It has how many flip-flops?
  \opts{5}{6}{7}{64}

  \item The minimum number of flip-flops for a state machine with 60 states is:
  \opts{5}{60}{6}{7}

  \item For the same sequence detector, which usually needs more states?
  \opts{Mealy}{Moore}{they always need the same}{it depends only on the clock frequency}

  \item How many transistors are in a static CMOS 3-input NAND?
  \opts{3}{4}{6}{8}

  \item PMOS transistors are usually drawn wider than NMOS transistors because:
  \opts{holes have lower mobility than electrons}{PMOS always has a higher threshold}{it is required for ESD protection}{it reduces leakage}

  \item NAND gates are generally preferred over NOR gates because:
  \opts{NAND uses fewer transistors}{NOR is not universal}{NAND has no parasitic capacitance}{NOR puts slow PMOS transistors in series}

  \item A 50\,$\Omega$ transmission line ends in a 25\,$\Omega$ load. The reflection coefficient is:
  \opts{$+1/3$}{$-1/3$}{$-1/2$}{0}

  \item Which best supplies a chip's power pins during a current step lasting a few nanoseconds?
  \opts{a large capacitor far from the chip}{no capacitor; the regulator responds instantly}{a small, low-inductance capacitor right at the pins}{an inductor in series at the pins}

  \item The main purpose of a heat sink's fins is to:
  \opts{store heat}{increase thermal resistance}{insulate the chip}{increase surface area for convection}

  \item A 300\,W chip has 0.2\,\textdegree C/W total thermal resistance to 30\,\textdegree C air. Its junction temperature is:
  \opts{60\,\textdegree C}{75\,\textdegree C}{90\,\textdegree C}{120\,\textdegree C}

  \item True or false: two boards exchanging digital signals must share a ground connection.

  \item In a 6T SRAM cell, read stability requires:
  \opts{access NMOS stronger than pull-down NMOS}{pull-up PMOS stronger than access NMOS}{pull-down NMOS stronger than access NMOS}{all transistors the same size}

  \item DRAM needs periodic refresh because:
  \opts{the stored charge leaks away}{reads are too slow}{neighboring SRAM interferes}{clock gating requires it}

  \item A block will be idle for a long time. What reduces its leakage power most?
  \opts{clock gating}{lowering only the frequency}{power gating}{adding decoupling capacitors}

  \item In the classic five-stage pipeline, loads read data memory in which stage?
  \opts{IF}{ID}{EX}{MEM}

  \item A cache has a 1\,ns hit time, a 10\% miss rate and a 50\,ns miss penalty. The AMAT is:
  \opts{5\,ns}{51\,ns}{6\,ns}{50\,ns}

  \item A two-flop synchronizer is appropriate for crossing:
  \opts{a 32-bit data bus}{a single-bit level signal}{the clock itself}{a multi-bit counter in binary}

  \item A writer at 100\,MHz sends bursts of 200 words; a reader at 50\,MHz reads one word per cycle. The minimum FIFO depth is:
  \opts{100}{50}{150}{200}

  \item An unbuffered wire's length doubles. Its delay becomes about:
  \opts{the same}{$2\times$}{$4\times$}{$8\times$}

  \item A divide-by-3 clock divider with a 50\% duty cycle requires:
  \opts{only rising clock edges}{both clock edges}{a PLL}{an asynchronous reset}

  \item Electromigration gets worse with:
  \opts{lower temperature}{lower supply voltage}{wider wires}{higher current density and temperature}

  \item Latch-up is prevented by:
  \opts{guard rings and frequent well and substrate taps}{lowering the clock frequency}{adding decoupling capacitors}{using NOR instead of NAND}

  \item BF16 has the same number of exponent bits as:
  \opts{FP16}{FP32}{INT8}{FP8 E4M3}

  \item In steady state, an ideal CMOS inverter draws:
  \opts{only leakage current, nearly zero}{a large current}{its saturation current}{a current that depends on the input}

  \item[\textbf{Challenge.}]\vspace{-0.5em}

  \item A chain of 50 XOR gates: gate 1 has inputs 1 and $X$; every later gate has inputs (previous gate's output) and $X$. The output of gate 50 is:
  \opts{$X$}{$\nt{X}$}{1}{0}

  \item Identical NMOS transistors, $\VDD = 5$\,V, $\Vt = 2$\,V. An input of 4\,V passes through one NMOS whose gate is at 5\,V. The output node reaches:
  \opts{4\,V}{3\,V}{2\,V}{5\,V}

  \item A 1\,MB, 4-way set-associative cache has 64-byte lines and 32-bit addresses. The tag is how many bits?
  \opts{12}{14}{16}{20}

  \item A 202-byte read starts at address \code{0x31121133}. Replies carry at most 64 bytes and may not cross 64-byte boundaries. How many reply packets?
  \opts{3}{4}{5}{6}

  \item A design moves from $(f, 2V)$ to $(4f, 3V)$. Dynamic power changes by a factor of:
  \opts{6}{9}{4}{12}

  \item Net A can switch between 1.0 and 1.6\,ns; adjacent net B between 1.4 and 2.2\,ns. Crosstalk between them is possible during:
  \opts{no time}{1.0 to 2.2\,ns}{1.4 to 1.6\,ns}{1.6 to 2.2\,ns}
\end{enumerate}

\section{Part B: What does this C code print?}

Assume a 64-bit machine.

\begin{enumerate}
  \item What does this print?
\begin{lstlisting}[language=C,numbers=none]
#define SQ(x) x*x
printf("%d", SQ(1+2));
\end{lstlisting}
  \item What does this print?
\begin{lstlisting}[language=C,numbers=none]
int f(void) { static int a = 12; return a++; }
int main(void) { f(); f(); printf("%d", f()); return 0; }
\end{lstlisting}
  \item What does this print?
\begin{lstlisting}[language=C,numbers=none]
int i = 8, j = 2;
printf("%d %d", i << j, i >> j);
\end{lstlisting}
  \item What does this print?
\begin{lstlisting}[language=C,numbers=none]
struct S { char c; void *p; };
printf("%zu", sizeof(struct S));
\end{lstlisting}
  \item What does this print?
\begin{lstlisting}[language=C,numbers=none]
unsigned char c = 250;
c += 10;
printf("%d", c);
\end{lstlisting}
\end{enumerate}

\section{Part C: Aptitude}

\begin{enumerate}
  \item How many trailing zeros does $50!$ have?
  \item Six ants on the corners of a hexagon each pick a random direction along the edges. What is the probability that none collide?
  \item One pipe fills a tank in 4 hours and another in 12 hours. How long do both together take?
  \item What is the maximum number of comparisons binary search needs on 4096 sorted elements?
  \item 1000 bulbs start off; person $k$ toggles every $k$th bulb, for $k = 1$ to 1000. How many bulbs end up on?
\end{enumerate}

\section{Part D: Coding}

\begin{enumerate}
  \item \textbf{Reverse a linked list.} Given the head of a singly linked list, reverse it and return the new head. Use $O(1)$ extra memory.
  \item \textbf{XOR subarrays.} Given an integer array and an integer $k$, return the number of contiguous subarrays whose elements XOR to $k$. The array can have up to $10^5$ elements.
\end{enumerate}

\section{Answer key}

\subsection*{Part A}

\begin{enumerate}
  \item \textbf{(c)} NMOS passes a weak 1: $\VDD - \Vtn = 0.7$\,V.
  \item \textbf{(b)} The pull-up network is PMOS, which passes a strong 1.
  \item \textbf{(d)} The top NMOS pulls up but stops one threshold below $\VDD$.
  \item \textbf{(c)} Input 1 gives a weak 1, input 0 gives a weak 0: non-inverting with degraded levels.
  \item \textbf{(d)} Edge triggering prevents race-through; (a), (b) and (c) are latch advantages.
  \item \textbf{(a)} $\nt{B}A + B = A + B$.
  \item \textbf{(b)} One mux makes $\nt{B}$; one selects $B$ or $\nt{B}$ by $A$.
  \item \textbf{(d)} $2 \times 2^2 = 8$.
  \item \textbf{(b)} $1.1^2 \times 1.2 = 1.452$.
  \item \textbf{(c)} More phonon scattering in both silicon and metal.
  \item \textbf{False.} A smaller gate raises the metal-to-gate area ratio, increasing the risk.
  \item \textbf{(a)} The diode drains plasma charge to the substrate.
  \item \textbf{(d)} Hold does not depend on the clock period.
  \item \textbf{(b)} $50 + 350 + 100 = 500$\,ps, so 2\,GHz.
  \item \textbf{(a)} The data gets extra time before the capture edge, but the hold window closes later.
  \item \textbf{(c)} A toggle flip-flop divides by 2.
  \item \textbf{(b)} $128/2 = 64 = 2^6$.
  \item \textbf{(c)} $\lceil\log_2 60\rceil = 6$.
  \item \textbf{(b)} Moore outputs are attached to states, which usually needs an extra state.
  \item \textbf{(c)} $2n$ for $n$ inputs.
  \item \textbf{(a)} Lower hole mobility is compensated by width.
  \item \textbf{(d)} Series PMOS is slow or must be very wide.
  \item \textbf{(b)} $(25-50)/(25+50) = -1/3$.
  \item \textbf{(c)} Fast steps are limited by inductance; small, close capacitors have the least.
  \item \textbf{(d)} Fins add area for the air to carry heat away.
  \item \textbf{(c)} $30 + 300 \times 0.2 = 90$\,\textdegree C.
  \item \textbf{True.} Voltages are only meaningful relative to a shared reference, and currents need a return path.
  \item \textbf{(c)} Otherwise the 0 node rises during a read and the cell can flip.
  \item \textbf{(a)} The capacitor leaks.
  \item \textbf{(c)} Power gating cuts the supply and thus most leakage; clock gating only stops switching.
  \item \textbf{(d)} MEM.
  \item \textbf{(c)} $1 + 0.1 \times 50 = 6$\,ns.
  \item \textbf{(b)} Buses and binary counters need Gray code, handshakes or FIFOs.
  \item \textbf{(a)} Burst lasts 2\,$\mu$s; 100 words are read in that time; $200 - 100 = 100$.
  \item \textbf{(c)} Delay grows with the square of length.
  \item \textbf{(b)} A falling-edge flip-flop provides the half-cycle delay.
  \item \textbf{(d)} Black's equation: worse with current density and exponentially with temperature.
  \item \textbf{(a)} Taps and guard rings keep the parasitic bipolar transistors off.
  \item \textbf{(b)} BF16 keeps FP32's 8-bit exponent.
  \item \textbf{(a)} One transistor is always off.
  \item[] \textbf{Challenge}
  \item \textbf{(c)} Odd gates output $\nt{X}$, even gates output 1.
  \item \textbf{(b)} $\min(4, 5 - 2) = 3$\,V.
  \item \textbf{(b)} Offset 6, index 12 (4096 sets), tag $32 - 18 = 14$.
  \item \textbf{(b)} 13 bytes to the first boundary, then 64, 64 and 61.
  \item \textbf{(b)} $4 \times (3/2)^2 = 9$.
  \item \textbf{(c)} The intersection of the two windows.
\end{enumerate}

\subsection*{Part B}

\begin{enumerate}
  \item \textbf{5.} \code{SQ(1+2)} expands to \code{1+2*1+2}.
  \item \textbf{14.} The calls return 12, 13, 14 (post-increment returns the old value; the static keeps counting).
  \item \textbf{32 2.}
  \item \textbf{16.} One byte, seven bytes of padding, an eight-byte pointer.
  \item \textbf{4.} $250 + 10 = 260$, which wraps modulo 256 to 4.
\end{enumerate}

\subsection*{Part C}

\begin{enumerate}
  \item \textbf{12.} $\lfloor 50/5\rfloor + \lfloor 50/25\rfloor = 10 + 2$.
  \item \textbf{1/32.} $2/2^6$.
  \item \textbf{3 hours.} Rates $\tfrac14 + \tfrac1{12} = \tfrac13$ per hour.
  \item \textbf{13.} $\lceil\log_2 4097\rceil = 13$.
  \item \textbf{31.} Perfect squares up to 1000: $1^2$ through $31^2 = 961$.
\end{enumerate}

\subsection*{Part D}

\begin{enumerate}
  \item The iterative \code{prev}/\code{cur}/\code{next} solution of Chapter~\ref{ch:prog}: $O(n)$ time, $O(1)$ space. Test it on an empty list, a single node and two nodes.
  \item Prefix XOR with a hash map counting earlier prefixes (Chapter~\ref{ch:prog}): for each prefix $P_j$, add the number of earlier prefixes equal to $P_j \oplus k$. Initialize the map with $\{0: 1\}$. $O(n)$ time; use a 64-bit counter because the answer can exceed $2^{31}$.
\end{enumerate}

\paragraph{Scoring yourself.} If you missed a question, do not just memorize the answer. Go back to the chapter that teaches it, reread the idea behind it, and redo the chapter's exercises. The real assessment will ask a different twist on the same idea.
